% Subsection: Boolean Combinations  (notes stub; content authored later)
\subsection{Boolean Combinations}

\begin{remark*}[Exposition]
Boolean combinations are the finite expressions obtained by using the logical
set operations together.  They are the operational core of algebras of sets:
start with finitely many subsets and build new subsets by repeatedly applying
union, intersection, and complement.
\end{remark*}

\begin{definition}[Finite Boolean Combination]\label{def:finite-boolean-combination}
Let \(X\) be a set and let \(A_1,\dots,A_n\subseteq X\), where \(n\geq 1\).
A \emph{finite Boolean combination of \(A_1,\dots,A_n\)} is any subset of
\(X\) obtained from \(A_1,\dots,A_n\) by finitely many applications of union,
intersection, and complement relative to \(X\).
\end{definition}

\begin{remark*}[Standard quantified statement]
If \(B\) is a finite Boolean combination of \(A_1,\dots,A_n\), then there is a
finite expression \(E\) built from the symbols
\[
  A_1,\dots,A_n,\quad \cup,\quad \cap,\quad {}^c
\]
such that \(B=E(A_1,\dots,A_n)\subseteq X\).
\end{remark*}

\begin{remark*}[Interpretation]
Finite Boolean combinations are the subsets describable by a finite logical
recipe using the input sets.  They are finite because the expression has only
finitely many operation symbols, even if the input sets themselves are
infinite.
\end{remark*}

\begin{remark*}[Examples]
For \(A,B,C\subseteq X\), each of the following is a finite Boolean
combination of \(A,B,C\):
\[
  A\cap B^c,
  \qquad
  (A\cup B)\cap C,
  \qquad
  (A^c\cap B)\cup(C\setminus A).
\]
Since \(C\setminus A=C\cap A^c\), set difference is expressible using
intersection and complement.
\end{remark*}

\begin{dependencies}
\begin{itemize}
  \item \hyperref[def:union]{Union}
  \item \hyperref[def:intersection]{Intersection}
  \item \hyperref[def:complement]{Complement}
  \item \hyperref[def:set-difference]{Set Difference}
  \item \hyperref[def:finite-infinite]{Finite and Infinite Sets}
\end{itemize}
\end{dependencies}
